\section{Numerical verification}
\label{app:numerics}

All numerical results quoted in the main text are produced by a
single unified script, \texttt{scripts/test\_crt\_decoupling.py},
which combines the three corpus tests \texttt{T1}, \texttt{T2}
and \texttt{T3} described below. The script uses
\texttt{mpmath} at $50$~decimal digits of precision for all
algebraic computations, and \texttt{numpy} with
\texttt{rng=np.random.default\_rng(seed=42)} for the empirical
Monte Carlo block of test~\texttt{T3}. The total runtime is under
five seconds on a 2024 laptop.

\subsection{T1: closed form for $\mathcal{I}_{p,q}$}
\label{appssec:T1}

\textbf{Block:} \texttt{T1} of \texttt{scripts/test\_crt\_decoupling.py}.

\textbf{Setup.} The block evaluates the closed
form~\eqref{eq:closed_form} at $k = 1, 2, \ldots, 7$, with
$\cos^2(\theta_3, \mustar) \approx 0.7808$ (couplage branch at
$\mustar = 15$) and $\sin^2(\theta_2, p_1) \approx 0.6914$
($\qplus(p_1=3) = 1/3$).

\textbf{Result.} The best fits are
\begin{align*}
  k_{3,5}^{\mathrm{eff}} &\;=\; 7,
  &
  \mathcal{I}^{\mathrm{pred}}_{3,5}
  &\;=\; 0.132 \text{ bits}
  &
  (\text{target } 0.138; \text{ error } 4.5\%),
  \\[2pt]
  k_{3,7}^{\mathrm{eff}} &\;=\; 4,
  &
  \mathcal{I}^{\mathrm{pred}}_{3,7}
  &\;=\; 0.306 \text{ bits}
  &
  (\text{target } 0.283; \text{ error } 8.1\%).
\end{align*}
Both fits lie within the $\pm 20\%$ tolerance of the corpus
T1~\cite[scripts/]{PT_Monograph}, hence
\textbf{T1 PASS} (with caveat: see Remark~\ref{rmk:k_interp}).
The full table at $k = 1, \ldots, 7$ is reproduced in
Table~\ref{tab:closed_form_check} of
Section~\ref{ssec:t1_validation}.

\subsection{T2: factorisation of the Ruelle measure}
\label{appssec:T2}

\textbf{Block:} \texttt{T2} of \texttt{scripts/test\_crt\_decoupling.py}.

\textbf{Setup.} The block builds the transfer matrices
$T_3 = \mathrm{antidiag}(1, 1)$ on $\mathcal{S}_3 = \{1, 2\}$ and
$T_5$ on $\mathcal{S}_5 = \{1, 2, 3, 4\}$ with off-diagonal
entries $1/(p-2) = 1/3$ as in Eq.~\eqref{eq:Tp_explicit}, computes
their left Perron eigenvectors by augmented least squares at
$50$~digits of precision, and compares the directly computed
$\pi_{15}^{\mathrm{Ruelle}}$ to the Kronecker product
$\pi_3^{\mathrm{Ruelle}} \otimes \pi_5^{\mathrm{Ruelle}}$.

\textbf{Result.}
\begin{itemize}[itemsep=2pt,leftmargin=*]
  \item Per-factor stationary distributions:
    $\pi_3^{\mathrm{Ruelle}} = (1/2, 1/2)$,
    $\pi_5^{\mathrm{Ruelle}} = (1/4, 1/4, 1/4, 1/4)$.
  \item Stationarity of $\pi_3 \otimes \pi_5$ under $T_{15}$:
    $\|(\pi_3 \otimes \pi_5) T_{15} - (\pi_3 \otimes \pi_5)\|_2
    = 0$ exactly at $50$~dps.
  \item Direct vs.\ tensor:
    $\|\pi_{15}^{\mathrm{direct}} - \pi_3 \otimes \pi_5\|_2 =
    5.88 \times 10^{-53}$ (machine precision at $50$~dps).
  \item Kullback--Leibler divergence:
    $\DKL(\pi_{15}^{\mathrm{direct}} \| \pi_3 \otimes \pi_5) =
    0$ to $50$~dps.
\end{itemize}

\textbf{T2 PASS.} This is the numerical witness of
Lemma~\ref{lem:pi_factor}.

\subsection{T3: empirical MI under the product measure}
\label{appssec:T3}

\textbf{Block:} \texttt{T3} of \texttt{scripts/test\_crt\_decoupling.py}.

\textbf{Setup.} The block draws $N = 10^6$ i.i.d.\ samples from
$\pi_3 \otimes \pi_5$ using the \texttt{numpy} random generator
with seed $42$, and estimates the empirical mutual
information of the two coordinate marginals by direct histogram
counts.

\textbf{Result.}
$\mathrm{MI}_{\mathrm{emp}} = 1.45 \times 10^{-6}$ bits,
to be compared with the finite-sample threshold
$1/\sqrt{N} \approx 10^{-3}$ bits. The observed value is three
orders of magnitude below the threshold, consistent with the
theoretical prediction of exactly zero mutual information under
the product measure (Theorem~\ref{thm:ruelle_zero}) plus Monte
Carlo noise.

\textbf{T3 PASS.}

\subsection{Reproducibility}
\label{appssec:reproducibility}

The script \texttt{scripts/test\_crt\_decoupling.py} prints all
three blocks T1, T2, T3 to standard output with PASS/FAIL flags
at the end of each block and a global \texttt{SYNTHESE} at the
end. Total wall-clock time on a 2024 laptop is under five
seconds, with the Monte Carlo block T3 dominating ($\sim 4$ s);
the algebraic blocks T1 and T2 each take well under one second
at $50$~dps precision.

To reproduce the values quoted above, run
\begin{quote}
  \texttt{python3 scripts/test\_crt\_decoupling.py}
\end{quote}
from the article root. The script has no command-line arguments
and requires only \texttt{mpmath} and \texttt{numpy} as
external dependencies.
